\documentclass[aspectratio=169]{beamer}
\usepackage{iftex}
\ifPDFTeX
  \usepackage[T1]{fontenc}
  \usepackage{lmodern}
\fi
\usepackage[czech]{babel}
\usepackage{amsmath}
\usepackage{graphicx}
\usepackage{xcolor}
\usepackage{listings}
\usepackage{tikz}
\usetikzlibrary{arrows.meta,positioning}

\newcommand{\CVUTTemplatePath}{../../../utils/presentation-template}
\usepackage{\CVUTTemplatePath/beamerthemeCVUT}
\setbeamertemplate{footline}[frame number]
\beamertemplatenavigationsymbolsempty

\newcommand{\LANG}{CZ}
\newcommand{\FRONTPAGE}{dark}

\title{Datové typy a vstup/výstup}
\subtitle{Základy programování\\Matouš Cejnek\\U12110, FS, ČVUT v Praze}
\author{}
\institute{}
\date{}
\renewcommand{\headlineA}{}
\renewcommand{\headlineB}{}
\renewcommand{\headlineC}{}

\definecolor{CodeBlue}{RGB}{38,83,125}
\definecolor{SoftGray}{RGB}{242,244,247}
\definecolor{BitGreen}{RGB}{38,125,94}
\setbeamerfont{normal text}{size=\large}
\setbeamerfont{frametitle}{size=\LARGE}
\setbeamercolor{block title}{fg=white,bg=CodeBlue}
\setbeamercolor{block body}{fg=black,bg=SoftGray}
\newcommand{\term}[1]{\textcolor{CodeBlue}{\textbf{#1}}}

\lstdefinestyle{python}{
  language=Python,
  basicstyle=\ttfamily\large,
  keywordstyle=\color{CodeBlue}\bfseries,
  stringstyle=\color{BitGreen},
  commentstyle=\color{gray},
  showstringspaces=false,
  columns=fullflexible,
  keepspaces=true
}

\begin{document}

% 0:00--0:01
\begin{frame}
  \clearpage\thispagestyle{empty}\titlepage
\end{frame}

% 0:01--0:03
\begin{frame}{Co si dnes odnesete?}
  \begin{itemize}
    \item Jak počítač ukládá informace pomocí bitů a bajtů?
    \item Jak se reprezentují celá a desetinná čísla?
    \item Jaké základní datové typy a kontejnery nabízí Python?
    \item Jak převádět hodnoty mezi typy?
    \item Jak program komunikuje s okolím přes vstup a výstup?
  \end{itemize}
\end{frame}

\begin{frame}{Na co navazujeme?}
  \begin{itemize}
    \item Program pracuje s \term{vstupem}, provádí algoritmus a vytváří \term{výstup}.
    \item Proměnná dovoluje hodnotu pojmenovat a později ji použít.
    \item Dnes se ptáme, \textbf{jaký význam hodnota má} a \textbf{jak je uložena}.
  \end{itemize}
  \begin{alertblock}{Dnes neopakujeme řízení toku}
    Podmínky a cykly procvičíme na navazujícím semináři a cvičení.
  \end{alertblock}
\end{frame}

% 0:03--0:11
\begin{frame}{Informace potřebuje reprezentaci}
  \begin{block}{Data}
    Informace zapsaná v podobě, kterou lze uložit, přenášet a zpracovat.
  \end{block}
  \begin{itemize}
    \item Stejný vzor bitů může podle dohody znamenat číslo, znak nebo barvu.
    \item \term{Datový typ} určuje možné hodnoty a povolené operace.
    \item Reprezentace propojuje abstraktní hodnotu s fyzickou pamětí.
  \end{itemize}
\end{frame}

\begin{frame}{Bit a bajt}
  \begin{itemize}
    \item \term{Bit} je jedna binární číslice: \texttt{0} nebo \texttt{1}.
    \item \term{Bajt (byte)} je běžně skupina osmi bitů.
    \item Osm bitů rozlišuje $2^8=256$ různých kombinací.
  \end{itemize}
  \vspace{5mm}
  \centering
  \Large\ttfamily 0\;1\;1\;0\;1\;0\;0\;1
  \par\normalsize\rmfamily
  \vspace{3mm}
  Jeden bajt; význam závisí na zvoleném typu a kódování.
\end{frame}

\begin{frame}{Binární zápis celého čísla}
  \begin{columns}[T]
    \begin{column}{0.48\textwidth}
      \centering
      \large\ttfamily 00101101\rmfamily
      \begin{align*}
        00101101_2
          &= 32+8+4+1 \\
          &= 45_{10}
      \end{align*}
    \end{column}
    \begin{column}{0.46\textwidth}
      \begin{itemize}
        \item Každá pozice má váhu mocniny dvou.
        \item Nejpravější bit má váhu $2^0=1$.
        \item Přidání bitu zdvojnásobí počet možných stavů.
      \end{itemize}
    \end{column}
  \end{columns}
\end{frame}

\begin{frame}{Celé číslo bez znaménka}
  \begin{block}{Unsigned integer}
    Všech $n$ bitů vyjadřuje velikost nezáporného celého čísla.
  \end{block}
  \[
    \text{rozsah: } 0 \text{ až } 2^n-1
  \]
  \begin{itemize}
    \item 8 bitů: 0 až 255
    \item 16 bitů: 0 až 65\,535
    \item Vhodné tam, kde záporná hodnota nedává smysl.
  \end{itemize}
\end{frame}

\begin{frame}{Jak uložit záporná čísla?}
  \begin{itemize}
    \item Potřebujeme do stejného počtu bitů zahrnout kladné i záporné hodnoty.
    \item Samostatný „bit znaménka“ je intuitivní, ale přináší dvě reprezentace nuly.
    \item Moderní počítače obvykle používají \term{dvojkový doplněk}.
  \end{itemize}
  \begin{block}{Důležitá myšlenka}
    Bity samy znaménko nemají. Znaménko vzniká až jejich interpretací podle dohodnuté reprezentace.
  \end{block}
\end{frame}

\begin{frame}{Celé číslo se znaménkem}
  \begin{block}{Signed integer ve dvojkovém doplňku}
    Nejvyšší bit má zápornou váhu; ostatní bity mají obvyklé kladné váhy.
  \end{block}
  \[
    \text{rozsah pro } n \text{ bitů: } -2^{n-1} \text{ až } 2^{n-1}-1
  \]
  \begin{itemize}
    \item 8 bitů: $-128$ až $127$
    \item Nula má jedinou reprezentaci.
    \item Stejné obvody mohou provádět sčítání kladných i záporných čísel.
  \end{itemize}
\end{frame}

\begin{frame}{Omezený počet bitů znamená omezený rozsah}
  \begin{itemize}
    \item Výsledek mimo rozsah pevného typu způsobí \term{přetečení} (overflow).
    \item V některých prostředích hodnota „přeteče“ na druhý konec rozsahu.
    \item Jinde program ohlásí chybu nebo použije větší reprezentaci.
  \end{itemize}
  \begin{alertblock}{Python je v tomto neobvyklý}
    Pythonový \texttt{int} se podle potřeby zvětšuje; limitem je dostupná paměť, ne pevně 32 nebo 64 bitů.
  \end{alertblock}
\end{frame}

% 0:11--0:20
\begin{frame}{Proč nestačí celá čísla?}
  \begin{itemize}
    \item Měření a výpočty často obsahují zlomkovou část.
    \item Potřebujeme reprezentovat velmi malá i velmi velká čísla.
    \item Pevný počet bitů ale stále dovoluje jen konečný počet hodnot.
  \end{itemize}
  \begin{block}{Kompromis}
    Čísla s plovoucí řádovou čárkou nabízejí široký rozsah za cenu omezené přesnosti.
  \end{block}
\end{frame}

\begin{frame}{Vědecký zápis jako inspirace}
  \[
    6{,}022\times10^{23}
  \]
  \begin{itemize}
    \item znaménko říká, zda je hodnota kladná nebo záporná,
    \item významné číslice určují přesnost,
    \item exponent posouvá řádovou čárku a určuje velikost.
  \end{itemize}
  \vspace{3mm}
  Počítačový \texttt{float} používá stejnou myšlenku, ale v soustavě o základu 2.
\end{frame}

\begin{frame}{IEEE 754: běžná reprezentace \texttt{float}}
  \centering
  \begin{tikzpicture}[node distance=0mm]
    \node[draw,fill=SoftGray,minimum width=1.3cm,minimum height=1.2cm] (s) {znaménko};
    \node[draw,fill=SoftGray,minimum width=2.7cm,minimum height=1.2cm,right=of s] (e) {exponent};
    \node[draw,fill=SoftGray,minimum width=5.3cm,minimum height=1.2cm,right=of e] (f) {zlomková část};
  \end{tikzpicture}
  \vspace{6mm}
  \begin{itemize}
    \item \term{binary32}: 1 + 8 + 23 bitů
    \item \term{binary64}: 1 + 11 + 52 bitů
    \item Pythonový \texttt{float} je na běžných implementacích binary64.
  \end{itemize}
\end{frame}

\begin{frame}{Co jednotlivé části znamenají?}
  Pro běžná normalizovaná čísla zjednodušeně:
  \[
    (-1)^{\text{znaménko}}\times 1{,}\text{zlomek}\times 2^{\text{posunutý exponent}}
  \]
  \begin{itemize}
    \item Znaménko volí kladnou nebo zápornou hodnotu.
    \item Exponent určuje řád neboli přibližnou velikost.
    \item Zlomková část nese významné binární číslice.
  \end{itemize}
  \begin{alertblock}{Pro tento kurz}
    Důležitější než ruční kódování bitů je rozumět omezenému rozsahu a přesnosti.
  \end{alertblock}
\end{frame}

\begin{frame}{Desetinné číslo nemusí být uloženo přesně}
  \begin{itemize}
    \item Konečný binární rozvoj mají například $0{,}5$ nebo $0{,}25$.
    \item Desetinná hodnota $0{,}1$ má v binární soustavě nekonečný periodický rozvoj.
    \item Uloží se proto nejbližší reprezentovatelná hodnota.
  \end{itemize}
  \begin{block}{Důsledek}
    Výsledky s \texttt{float} porovnáváme s vhodnou tolerancí, ne vždy přesnou rovností.
  \end{block}
\end{frame}

\begin{frame}[fragile]{Zaokrouhlovací chyba není chyba Pythonu}
\begin{lstlisting}[style=python]
result = 0.1 + 0.2
print(result)          # 0.30000000000000004
\end{lstlisting}
  \begin{itemize}
    \item Operace začíná už mírně zaokrouhlenými vstupy.
    \item Další operace mohou chybu zvětšit nebo naopak skrýt.
    \item Pro peníze a přesnou desítkovou aritmetiku existují jiné typy, například \texttt{Decimal}.
  \end{itemize}
\end{frame}

\begin{frame}{Zvláštní hodnoty IEEE 754}
  \begin{itemize}
    \item $+\infty$ a $-\infty$: například výsledek překročení rozsahu v některých výpočtech,
    \item \term{NaN} (Not a Number): neurčitý nebo nedefinovaný výsledek,
    \item $+0$ a $-0$: dvě znaménkové podoby nuly,
    \item subnormální čísla: velmi malé hodnoty blízko nule.
  \end{itemize}
  \begin{alertblock}{NaN je zvláštní}
    NaN se nerovná ani samo sobě; ke kontrole používáme určené funkce.
  \end{alertblock}
\end{frame}

% 0:20--0:29
\begin{frame}{Základní skalární typy v Pythonu}
  \centering
  \begin{tabular}{lll}
    \textbf{Typ} & \textbf{Význam} & \textbf{Příklad} \\
    \hline
    \texttt{int} & celé číslo & \texttt{-42} \\
    \texttt{float} & číslo s plovoucí čárkou & \texttt{3.14} \\
    \texttt{bool} & pravdivostní hodnota & \texttt{True} \\
    \texttt{str} & text & \texttt{"ahoj"} \\
    \texttt{NoneType} & nepřítomnost hodnoty & \texttt{None} \\
  \end{tabular}
  \vspace{5mm}

  Typ můžeme zjistit funkcí \texttt{type(value)}.
\end{frame}

\begin{frame}{\texttt{bool} a \texttt{None} nejsou totéž}
  \begin{itemize}
    \item \texttt{bool} má dvě hodnoty: \texttt{True} a \texttt{False}.
    \item \texttt{None} vyjadřuje, že hodnota není k dispozici nebo nebyla určena.
    \item \texttt{0}, prázdný text a \texttt{None} mají odlišný význam.
  \end{itemize}
  \begin{block}{Typ zachycuje význam}
    „Naměřili jsme nulu“ není totéž jako „měření chybí“.
  \end{block}
\end{frame}

\begin{frame}{Text: znaky nejsou bajty}
  \begin{itemize}
    \item Pythonový \texttt{str} je posloupnost znaků Unicode.
    \item \term{Unicode} přiřazuje znakům kódové body: písmena, symboly i emoji.
    \item Při uložení nebo přenosu se text převádí na bajty pomocí \term{kódování}.
    \item Běžné kódování je UTF-8; jeden znak v něm může zabrat různý počet bajtů.
  \end{itemize}
  \begin{alertblock}{Častý omyl}
    Počet znaků textu nemusí být stejný jako počet bajtů v souboru.
  \end{alertblock}
\end{frame}

\begin{frame}{Proměnná, hodnota a typ}
  \begin{itemize}
    \item Proměnná je \term{jméno}, které odkazuje na objekt v paměti.
    \item Typ patří hodnotě (objektu), ne natrvalo jménu proměnné.
    \item Python určuje typ za běhu a jedno jméno může později odkazovat na jiný typ.
  \end{itemize}
  \begin{block}{Dynamické typování}
    Nemusíme typ proměnné předem deklarovat, ale typy hodnot stále určují povolené operace.
  \end{block}
\end{frame}

\begin{frame}[fragile]{Převod typu musí být výslovný}
\begin{lstlisting}[style=python]
count_text = "12"
count = int(count_text)
average = float("3.5")
label = str(42)
\end{lstlisting}
  \begin{itemize}
    \item \texttt{int()}, \texttt{float()} a \texttt{str()} vytvářejí hodnotu požadovaného typu.
    \item Ne každý převod dává smysl: \texttt{int("ahoj")} skončí chybou.
    \item Převod může ztrácet informaci: \texttt{int(3.9)} je \texttt{3}.
  \end{itemize}
\end{frame}

% 0:29--0:37
\begin{frame}{Kontejner uchovává více hodnot}
  \begin{block}{Kontejnerový typ}
    Hodnota, která sdružuje jiné hodnoty a určuje způsob přístupu k nim.
  \end{block}
  \begin{itemize}
    \item Záleží na pořadí prvků?
    \item Mohou se prvky opakovat?
    \item Lze kontejner po vytvoření měnit?
    \item Přistupujeme pozicí, hodnotou, nebo klíčem?
  \end{itemize}
\end{frame}

\begin{frame}{Čtyři běžné kontejnery Pythonu}
  \centering
  \small
  \begin{tabular}{llll}
    \textbf{Typ} & \textbf{Uspořádání} & \textbf{Měnitelný} & \textbf{Typické použití} \\
    \hline
    \texttt{list} & ano & ano & posloupnost prvků \\
    \texttt{tuple} & ano & ne & pevná skupina hodnot \\
    \texttt{set} & bez pozičního pořadí & ano & unikátní prvky \\
    \texttt{dict} & pořadí vložení & ano & dvojice klíč--hodnota \\
  \end{tabular}
  \par
  \normalsize
  \vspace{5mm}
  Volba kontejneru vyjadřuje, jak s daty zamýšlíme pracovat.
\end{frame}

\begin{frame}{\texttt{list} a \texttt{tuple}: uspořádané hodnoty}
  \begin{columns}[T]
    \begin{column}{0.46\textwidth}
      \begin{block}{\texttt{list}}
        \texttt{["red", "green", "blue"]}
      \end{block}
      \begin{itemize}
        \item lze měnit,
        \item přístup pomocí indexu,
        \item duplicity jsou povolené.
      \end{itemize}
    \end{column}
    \begin{column}{0.46\textwidth}
      \begin{block}{\texttt{tuple}}
        \texttt{(1920, 1080)}
      \end{block}
      \begin{itemize}
        \item po vytvoření se nemění,
        \item přístup pomocí indexu,
        \item vhodná pevná struktura.
      \end{itemize}
    \end{column}
  \end{columns}
\end{frame}

\begin{frame}{\texttt{set}: množina unikátních hodnot}
  \begin{itemize}
    \item Prvek je v množině nejvýše jednou.
    \item Nepracujeme s jeho pozičním indexem.
    \item Množina přirozeně podporuje sjednocení, průnik a rozdíl.
    \item Typická otázka zní: „Je tato hodnota přítomná?“
  \end{itemize}
  \begin{block}{Příklad významu}
    \texttt{allowed\_roles = \{"student", "teacher"\}}
  \end{block}
\end{frame}

\begin{frame}{\texttt{dict}: hodnoty pojmenované klíči}
  \begin{itemize}
    \item Slovník ukládá dvojice \term{klíč--hodnota}.
    \item Každý klíč je ve slovníku jedinečný.
    \item Hodnotu hledáme podle klíče, nikoli podle číselné pozice.
  \end{itemize}
  \begin{block}{Příklad významu}
    \texttt{person = \{"name": "Eva", "age": 21\}}
  \end{block}
  Slovník se hodí pro záznam s pojmenovanými položkami nebo pro rychlé vyhledávání.
\end{frame}

\begin{frame}{Měnitelnost ovlivňuje chování programu}
  \begin{itemize}
    \item \term{Mutable}: objekt lze změnit na místě --- například \texttt{list}, \texttt{set}, \texttt{dict}.
    \item \term{Immutable}: po vytvoření se nemění --- například \texttt{int}, \texttt{float}, \texttt{str}, \texttt{tuple}.
    \item Dvě proměnné mohou odkazovat na tentýž měnitelný objekt.
  \end{itemize}
  \begin{alertblock}{Praktický důsledek}
    Změna sdíleného seznamu může být viditelná přes více různých jmen.
  \end{alertblock}
\end{frame}

% 0:37--0:44
\begin{frame}{Program komunikuje pomocí proudů}
  \centering
  \begin{tikzpicture}[node distance=13mm,>=Latex]
    \node[draw,rounded corners,fill=SoftGray,minimum width=2.3cm] (in) {vstup};
    \node[draw,rounded corners,fill=SoftGray,minimum width=3.2cm,right=of in] (program) {program};
    \node[draw,rounded corners,fill=SoftGray,minimum width=2.3cm,right=of program] (out) {výstup};
    \draw[->,very thick] (in)--node[above]{data}(program);
    \draw[->,very thick] (program)--node[above]{data}(out);
  \end{tikzpicture}
  \vspace{7mm}
  \begin{itemize}
    \item \term{Proud (stream)} je postupná posloupnost dat.
    \item Zdroj nebo cíl může být terminál, soubor, zařízení i jiný program.
    \item Program nemusí znát fyzické provedení druhého konce proudu.
  \end{itemize}
\end{frame}

\begin{frame}{Standardní vstup a výstupy}
  \begin{itemize}
    \item \term{stdin}: standardní vstup, obvykle klávesnice nebo přesměrovaná data,
    \item \term{stdout}: běžný výstup programu,
    \item \term{stderr}: diagnostické a chybové zprávy.
  \end{itemize}
  \begin{block}{Proč jsou oddělené?}
    Výsledek programu lze uložit nebo předat dál, zatímco chyby zůstanou viditelné uživateli.
  \end{block}
\end{frame}

\begin{frame}[fragile]{Výstup v Pythonu: \texttt{print()}}
\begin{lstlisting}[style=python]
name = "Eva"
score = 42
print("Student:", name, "body:", score)
\end{lstlisting}
  \begin{itemize}
    \item \texttt{print()} převede hodnoty na text a zapíše je na standardní výstup.
    \item Více argumentů standardně oddělí mezerou.
    \item Na konec standardně přidá znak nového řádku.
  \end{itemize}
\end{frame}

\begin{frame}[fragile]{Vstup v Pythonu: \texttt{input()}}
\begin{lstlisting}[style=python]
age_text = input("Zadejte vek: ")
age = int(age_text)
\end{lstlisting}
  \begin{itemize}
    \item \texttt{input()} zobrazí výzvu a čeká na řádek standardního vstupu.
    \item Výsledkem je vždy \texttt{str}, i když uživatel napsal číslice.
    \item Význam vstupu určíme až kontrolou a případným převodem typu.
  \end{itemize}
\end{frame}

\begin{frame}{Rozhraní programu je dohoda}
  \begin{itemize}
    \item Program musí vědět, v jakém formátu vstup očekává.
    \item Uživatel nebo jiný program musí vědět, co výstup znamená.
    \item Typ, jednotka, pořadí i kódování jsou součástí této dohody.
  \end{itemize}
  \begin{alertblock}{„Zadejte číslo“ často nestačí}
    Celé, nebo desetinné? V jakém rozsahu? S desetinnou čárkou, nebo tečkou? V jaké jednotce?
  \end{alertblock}
\end{frame}

% 0:44--0:45
\begin{frame}{Co si zapamatovat?}
  \begin{itemize}
    \item Bity získávají význam teprve díky zvolené reprezentaci a datovému typu.
    \item Signed a unsigned celá čísla rozdělují stejný počet bitů jinak.
    \item IEEE 754 nabízí velký rozsah, ale mnoho hodnot ukládá pouze přibližně.
    \item Python má skalární typy i kontejnery s různými vlastnostmi.
    \item \texttt{input()} vrací text; převod na číslo je samostatný krok.
    \item Standardní proudy oddělují vstup, běžný výstup a chybové zprávy.
  \end{itemize}
  \vspace{3mm}
  \centering\Large Otázky?
\end{frame}

\end{document}
